\documentclass[11pt,a4paper]{article}

% ── Codificación, fuentes y lenguaje ───────────────────────────────────────────
\usepackage[utf8]{inputenc}
\usepackage[T1]{fontenc}
\usepackage[default,scale=0.95]{sourcesanspro}
\renewcommand{\familydefault}{\sfdefault}
\usepackage[spanish,es-noshorthands,es-tabla]{babel}

% ── Geometría y espaciado ─────────────────────────────────────────────────────
\usepackage[top=2.2cm, bottom=2.5cm, left=2.2cm, right=2.2cm, headheight=15pt]{geometry}
\usepackage{setspace}
\setstretch{1.18}
\usepackage{parskip}

% ── Matemáticas y unidades ────────────────────────────────────────────────────
\usepackage{amsmath}
\usepackage{amssymb}
\usepackage{siunitx}

% ── Tablas ────────────────────────────────────────────────────────────────────
\usepackage{booktabs}
\usepackage{tabularx}
\usepackage{array}
\usepackage{longtable}
\usepackage{colortbl}
\usepackage{enumitem}

% ── Colores, cajas e iconos ───────────────────────────────────────────────────
\usepackage[dvipsnames]{xcolor}
\usepackage{tcolorbox}
\tcbuselibrary{skins, breakable}
\usepackage{fontawesome5}
\usepackage{graphicx}

% ── Listados de código (dark-mode infográfico) ────────────────────────────────
\usepackage{listings}

% ── Secciones con barra de color (infografía) ────────────────────────────────
\usepackage{titlesec}
\titleformat{\section}
  {\normalfont\LARGE\bfseries\color{azulNoche}}
  {\colorbox{cyanNeon}{\color{fondoOscuro}\thesection}}{0.7em}{}
  [\vspace{2pt}\color{cyanNeon}\rule{\linewidth}{1.8pt}]
\titlespacing*{\section}{0pt}{24pt}{10pt}
\titleformat{\subsection}
  {\normalfont\large\bfseries\color{azulNoche}}
  {\textcolor{cyanNeon}{\thesubsection}}{0.7em}{}
  [\vspace{1pt}\color{grisLinea}\rule{\linewidth}{0.6pt}]
\titlespacing*{\subsection}{0pt}{14pt}{6pt}
\titleformat{\subsubsection}
  {\normalfont\normalsize\bfseries\color{azulNoche}}
  {\textcolor{cyanNeon}{\thesubsubsection}}{0.7em}{}
\titlespacing*{\subsubsection}{0pt}{10pt}{4pt}

% ── Cabeceras y pies de página ────────────────────────────────────────────────
\usepackage{fancyhdr}
\usepackage{lastpage}
\pagestyle{fancy}
\fancyhf{}
\renewcommand{\headrulewidth}{0pt}
\renewcommand{\footrulewidth}{0pt}
\fancyfoot[C]{%
  \begin{minipage}{\linewidth}
    \vspace{2pt}
    \color{cyanNeon}\rule{\linewidth}{1.2pt}\\[2pt]
    \centering\color{grisTexto}\footnotesize
    Página \thepage\ de \pageref{LastPage}
  \end{minipage}%
}

% ── Hiperenlaces ──────────────────────────────────────────────────────────────
\usepackage[colorlinks=true, linkcolor=azulNoche, urlcolor=cyanNeon]{hyperref}
\sloppy
\emergencystretch 3em

% ── Paleta de colores — tecnológico dark-mode ──────────────────────────────────
\definecolor{fondoOscuro}{HTML}{0D1B2A}
\definecolor{verdeTurquesa}{HTML}{004D40}
\definecolor{azulNoche}{HTML}{1B3A6B}
\definecolor{azulMedio}{HTML}{2A5298}
\definecolor{cyanNeon}{HTML}{00C8E0}
\definecolor{verdeSignal}{HTML}{00B887}
\definecolor{naranjaVivo}{HTML}{FF7043}
\definecolor{rojoAlerta}{HTML}{E53935}
\definecolor{grisPapel}{HTML}{F4F6FA}
\definecolor{grisLinea}{HTML}{C8D0E0}
\definecolor{grisTexto}{HTML}{6B7A99}
\definecolor{amarilloNota}{HTML}{FFB300}
\definecolor{codigoFondo}{HTML}{1E2D3D}
\definecolor{codigoComentario}{HTML}{7A8FA6}
\definecolor{codigoCadena}{HTML}{FF9D5C}
\definecolor{codigoKeyword}{HTML}{00C8E0}
\definecolor{codigoNumero}{HTML}{00E5A0}

\newcommand{\iconoBanda}{microchip}
\newcommand{\bandaTitulo}[3][fondoOscuro]{%
  \begin{tcolorbox}[
    enhanced, colback=#1, colframe=#1, boxrule=0pt, arc=8pt, outer arc=8pt,
    width=\linewidth, top=18pt, bottom=6pt, left=14pt, right=14pt,
    borderline south={3.5pt}{0pt}{cyanNeon},
  ]
    \begin{minipage}[c]{0.07\linewidth}
      \centering
      \textcolor{cyanNeon}{\fontsize{30}{30}\selectfont\faIcon{\iconoBanda}}
    \end{minipage}%
    \hspace{8pt}%
    \begin{minipage}[c]{0.88\linewidth}
      {\fontsize{20}{24}\selectfont\bfseries\color{white}#2}\par\vspace{5pt}%
      \textcolor{cyanNeon!80!white}{\small\faIcon{angle-right}~#3}%
    \end{minipage}
    \vspace{4pt}
  \end{tcolorbox}%
  \vspace{8pt}%
}

\tcbset{
  tarjetaDato/.style={enhanced, breakable, arc=6pt, outer arc=6pt, colback=grisPapel, colframe=azulNoche, boxrule=1pt, fonttitle=\bfseries\small, coltitle=white, attach boxed title to top left={yshift=-3mm, xshift=6mm}, boxed title style={colback=azulNoche, colframe=azulNoche, arc=4pt, boxrule=0pt}, drop fuzzy shadow=azulNoche!30!white, top=8pt, bottom=8pt, left=10pt, right=10pt},
  tarjetaIcono/.style={enhanced, breakable, arc=6pt, outer arc=6pt, colback=white, colframe=grisLinea, boxrule=0.8pt, fonttitle=\bfseries\small, coltitle=azulNoche, borderline west={3pt}{0pt}{cyanNeon}, drop fuzzy shadow=azulNoche!25!white, top=6pt, bottom=6pt, left=10pt, right=10pt},
  cajaContenido/.style={enhanced, breakable, arc=5pt, outer arc=5pt, colback=grisPapel, colframe=grisLinea, boxrule=0.8pt, fonttitle=\bfseries\small, coltitle=azulNoche, top=6pt, bottom=6pt, left=10pt, right=10pt},
  cajaFortaleza/.style={enhanced, breakable, arc=6pt, outer arc=6pt, colback=verdeSignal!8!white, colframe=verdeSignal, boxrule=1pt, borderline west={4pt}{0pt}{verdeSignal}, fonttitle=\bfseries\small, coltitle=verdeSignal!70!black, drop fuzzy shadow=verdeSignal!20!white, top=6pt, bottom=6pt, left=10pt, right=10pt},
  cajaMejora/.style={enhanced, breakable, arc=6pt, outer arc=6pt, colback=naranjaVivo!8!white, colframe=naranjaVivo, boxrule=1pt, borderline west={4pt}{0pt}{naranjaVivo}, fonttitle=\bfseries\small, coltitle=naranjaVivo!80!black, drop fuzzy shadow=naranjaVivo!20!white, top=6pt, bottom=6pt, left=10pt, right=10pt},
  cajaRecomendacion/.style={enhanced, breakable, arc=6pt, outer arc=6pt, colback=cyanNeon!6!white, colframe=cyanNeon!60!azulNoche, boxrule=1pt, borderline west={4pt}{0pt}{cyanNeon}, fonttitle=\bfseries\small, coltitle=azulNoche, drop fuzzy shadow=azulNoche!20!white, top=6pt, bottom=6pt, left=10pt, right=10pt},
}

\newcommand{\iconotexto}[2]{\textcolor{cyanNeon}{\faIcon{#1}}~#2}

\lstdefinestyle{estiloCodigo}{
    backgroundcolor=\color{codigoFondo}, basicstyle=\ttfamily\small\color{white},
    commentstyle=\color{codigoComentario}\itshape, stringstyle=\color{codigoCadena},
    keywordstyle=\color{codigoKeyword}\bfseries, numberstyle=\tiny\color{codigoComentario},
    numbers=left, numbersep=10pt, showstringspaces=false, breaklines=true, breakatwhitespace=false,
    tabsize=4, frame=none, captionpos=b, aboveskip=6pt, belowskip=4pt, xleftmargin=14pt, framexleftmargin=14pt,
    literate={á}{{\'a}}1 {é}{{\'e}}1 {í}{{\'i}}1 {ó}{{\'o}}1 {ú}{{\'u}}1 {Á}{{\'A}}1 {É}{{\'E}}1 {Í}{{\'I}}1 {Ó}{{\'O}}1 {Ú}{{\'U}}1 {ñ}{{\~n}}1 {Ñ}{{\~N}}1 {ü}{{\"u}}1 {Ü}{{\"U}}1 {¿}{{?`}}1 {¡}{{!`}}1
}
\lstset{style=estiloCodigo}
\BeforeBeginEnvironment{lstlisting}{%
  \begin{tcolorbox}[enhanced, breakable, arc=6pt, outer arc=6pt, colback=codigoFondo, colframe=azulNoche, boxrule=1pt, drop fuzzy shadow=azulNoche!25!white, top=2pt, bottom=2pt, left=0pt, right=0pt, colupper=white, coltext=white]%
}
\AfterEndEnvironment{lstlisting}{\end{tcolorbox}}

\newtcolorbox{cajaRecuerda}{enhanced, breakable, arc=6pt, outer arc=6pt, colback=azulNoche!10!grisPapel, colframe=azulNoche, boxrule=1pt, borderline west={4pt}{0pt}{cyanNeon}, fonttitle=\bfseries\small\color{white}, title={\faIcon{info-circle}~Recuerda}, attach boxed title to top left={yshift=-3mm, xshift=6mm}, boxed title style={colback=azulNoche, arc=4pt, boxrule=0pt}, drop fuzzy shadow=azulNoche!25!white, left=10pt, right=10pt, top=8pt, bottom=8pt}
\newtcolorbox{cajaConcepto}{enhanced, breakable, arc=6pt, outer arc=6pt, colback=verdeSignal!8!white, colframe=verdeSignal, boxrule=1pt, borderline west={4pt}{0pt}{verdeSignal}, fonttitle=\bfseries\small\color{white}, title={\faIcon{lightbulb}~Concepto clave}, attach boxed title to top left={yshift=-3mm, xshift=6mm}, boxed title style={colback=verdeSignal!80!black, arc=4pt, boxrule=0pt}, drop fuzzy shadow=verdeSignal!20!white, left=10pt, right=10pt, top=8pt, bottom=8pt}
\newtcolorbox{cajaEjemplo}{enhanced, breakable, arc=6pt, outer arc=6pt, colback=cyanNeon!5!white, colframe=cyanNeon!70!azulNoche, boxrule=1pt, borderline west={4pt}{0pt}{cyanNeon}, fonttitle=\bfseries\small\color{fondoOscuro}, title={\faIcon{code}~Ejemplo}, attach boxed title to top left={yshift=-3mm, xshift=6mm}, boxed title style={colback=cyanNeon, arc=4pt, boxrule=0pt}, drop fuzzy shadow=azulNoche!20!white, left=10pt, right=10pt, top=8pt, bottom=8pt}
\newtcolorbox{cajaEjercicio}{enhanced, breakable, arc=6pt, outer arc=6pt, colback=azulMedio!7!white, colframe=azulMedio, boxrule=1pt, borderline west={4pt}{0pt}{azulMedio}, fonttitle=\bfseries\small\color{white}, title={\faIcon{pencil-alt}~Ejercicio}, attach boxed title to top left={yshift=-3mm, xshift=6mm}, boxed title style={colback=azulMedio, arc=4pt, boxrule=0pt}, drop fuzzy shadow=azulNoche!20!white, left=10pt, right=10pt, top=8pt, bottom=8pt}
\newtcolorbox{cajaReto}{enhanced, breakable, arc=6pt, outer arc=6pt, colback=naranjaVivo!6!white, colframe=naranjaVivo, boxrule=1pt, borderline west={4pt}{0pt}{naranjaVivo}, fonttitle=\bfseries\small\color{white}, title={\faIcon{fire}~Reto}, attach boxed title to top left={yshift=-3mm, xshift=6mm}, boxed title style={colback=naranjaVivo, arc=4pt, boxrule=0pt}, drop fuzzy shadow=naranjaVivo!20!white, left=10pt, right=10pt, top=8pt, bottom=8pt}

% ─────────────────────────────────────────────────────────────────────────────
\renewcommand{\iconoBanda}{atom}
\fancyhead[L]{\footnotesize\color{grisTexto}\faIcon{atom}~LABORATORIO\_II\_FISICA}
\fancyhead[R]{\footnotesize\color{grisTexto}Guía de repositorio GitHub}
\hypersetup{pdftitle={Guía para crear y publicar tu repositorio en GitHub},pdfauthor={Prof. César Rodríguez},pdfsubject={LABORATORIO_II_FISICA}}

\begin{document}
\bandaTitulo[fondoOscuro]{Guía para crear y publicar tu repositorio en GitHub}{Laboratorio II de Física --- estructura de carpetas, plantilla y primer push}

\begin{center}\small\color{grisTexto}Prof. César Rodríguez \textbar{} Universidad de Oriente, Núcleo de Sucre, Departamento de Física \textbar{} 24 de noviembre de 2026\end{center}

\begin{tcolorbox}[tarjetaDato,title={\faIcon{bullseye}~Objetivo de esta guía}]
Esta guía te lleva, paso a paso, desde una carpeta vacía en tu computadora hasta un repositorio privado en GitHub con la estructura oficial del curso \textbf{Laboratorio II de Física}. Al terminar tendrás la plantilla del grupo lista para trabajar durante todo el semestre y sabrás cómo subir cambios con Git.
\end{tcolorbox}

\section{\iconotexto{toolbox}{Qué necesitas antes de empezar}}

\begin{tcolorbox}[tarjetaIcono,title={\faIcon{check-circle}~Requisitos mínimos}]
\begin{itemize}[leftmargin=1.4em,itemsep=3pt,topsep=2pt]
  \item \textbf{Cuenta de GitHub.} Si no tienes una, créala en \url{https://github.com/signup} con tu correo.
  \item \textbf{Git instalado en tu PC.} Compruébalo en una terminal:
\end{itemize}
\begin{lstlisting}
git --version
\end{lstlisting}
Si ves un número de versión (por ejemplo \texttt{git version 2.x.x}), Git está listo. Si no, instálalo desde \url{https://git-scm.com/downloads} y reinicia la terminal.
\begin{tcolorbox}[cajaContenido]
\textbf{Identidad de Git (solo la primera vez en esa PC):}
\begin{lstlisting}
git config --global user.name "Tu Nombre"
git config --global user.email "tu-correo@ejemplo.com"
\end{lstlisting}
Usa el mismo correo asociado a tu cuenta de GitHub para que los commits queden vinculados a tu perfil.
\end{tcolorbox}
\end{tcolorbox}

\begin{cajaRecuerda}
Necesitarás una terminal: \textbf{Linux/macOS} usan Terminal; \textbf{Windows} puede usar Git Bash (se instala junto con Git), PowerShell o WSL. Los comandos de esta guía funcionan igual en Git Bash/Linux/macOS. En PowerShell, \texttt{mkdir -p} puede requerir \texttt{New-Item -ItemType Directory -Force}.
\end{cajaRecuerda}

\section{\iconotexto{sitemap}{Conoce la plantilla oficial del curso}}

La plantilla oficial define la estructura que el docente espera ver en tu repositorio. No inventes otra estructura: parte de esta plantilla y personalízala.

\begin{tcolorbox}[tarjetaDato,title={\faIcon{github}~Ubicación de la plantilla}]
\begin{itemize}[leftmargin=1.4em,itemsep=4pt,topsep=2pt]
  \item \textbf{Curso:} Laboratorio II de Física.
  \item \textbf{Carpeta del curso en el repositorio de referencia:}\\
  \url{https://github.com/cero97-ctrl/ELECTRONICA/tree/master/cursos/LABORATORIO_II_FISICA}
  \item \textbf{Plantilla del grupo:}\\
  \url{https://github.com/cero97-ctrl/ELECTRONICA/tree/master/cursos/LABORATORIO_II_FISICA/plantillas/repo_grupo}
\end{itemize}
\end{tcolorbox}

\subsection{\iconotexto{folder-open}{Estructura obligatoria}}

Cada semana debe tener esta forma. En Lab II, las 10 semanas corresponden a prácticas de laboratorio eléctrico/electrónico; la estructura de carpetas es la misma para todas las semanas.

\begin{lstlisting}
repo_grupo/
|-- README.md                  # integrantes, roles, tabla de semanas
|-- .gitignore                 # secretos, venv, auxiliares LaTeX, etc.
|-- semana_01/
|   |-- informe.pdf            # informe de la practica/taller
|   |-- datos/                 # CSV crudos, sin editar
|   |-- analisis/              # scripts que reproducen figuras/resultados
|   |-- figuras/               # graficas generadas por los scripts
|   `-- README.md              # que se hizo y que no se pudo verificar
|-- semana_02/
|   `-- ...
...
`-- semana_10/
    |-- informe.pdf
    |-- datos/
    |-- analisis/
    |-- figuras/
    |-- portafolio/            # cierre del semestre
    `-- README.md
\end{lstlisting}

\begin{tcolorbox}[cajaRecomendacion,title={\faIcon{info-circle}~Regla de oro}]
La entrega es el \textbf{repositorio completo}, no un archivo suelto. Los datos crudos no se editan; toda corrección va en \texttt{analisis/}. Cada cifra del informe debe salir de un CSV o de un script entregado.
\end{tcolorbox}

\section{\iconotexto{laptop-code}{Crear la estructura en tu PC}}

Elige \textbf{una} de las dos opciones. La opción A es la recomendada porque evita errores de tipeo.

\subsection{\iconotexto{copy}{Opción A (recomendada): copiar la plantilla oficial}}

\begin{enumerate}[leftmargin=1.6em,itemsep=4pt,topsep=2pt]
  \item Abre la URL de la plantilla: \url{https://github.com/cero97-ctrl/ELECTRONICA/tree/master/cursos/LABORATORIO_II_FISICA/plantillas/repo_grupo}.
  \item Descarga la carpeta \texttt{plantillas/repo\_grupo} del curso. Puedes hacerlo clonando el repositorio de referencia o descargando los archivos visibles en GitHub. Si clonas, hazlo en una carpeta temporal, \textbf{no} dentro de tu repo final.
  \item Copia el contenido de \texttt{repo\_grupo} a una carpeta nueva con el nombre de tu grupo, por ejemplo:
\end{enumerate}

\begin{lstlisting}
# Ejemplo en Linux/macOS/Git Bash
mkdir -p ~/laboratorios
cp -r /ruta/a/repo_grupo ~/laboratorios/lab2-fisica-grupo-01
cd ~/laboratorios/lab2-fisica-grupo-01
\end{lstlisting}

\begin{cajaRecuerda}
Copia el \textbf{contenido} de la plantilla (incluido el archivo oculto \texttt{.gitignore}). En Linux/macOS, \texttt{cp -r repo\_grupo/. destino/} copia también archivos ocultos. Verifica con \texttt{ls -a}.
\end{cajaRecuerda}

\subsection{\iconotexto{terminal}{Opción B: crear la estructura desde cero}}

Si prefieres crearla manualmente, ejecuta estos comandos en la carpeta donde quieras trabajar. Ajusta el nombre \texttt{lab2-fisica-grupo-01}.

\begin{lstlisting}
mkdir -p lab2-fisica-grupo-01
cd lab2-fisica-grupo-01

# Crear semanas con sus subcarpetas
for i in 01 02 03 04 05 06 07 08 09 10; do
  mkdir -p "semana_$i/datos" "semana_$i/analisis" "semana_$i/figuras"
  touch "semana_$i/datos/.gitkeep"
  touch "semana_$i/analisis/.gitkeep"
  touch "semana_$i/figuras/.gitkeep"
  touch "semana_$i/README.md"
done
mkdir -p "semana_10/portafolio"
touch "semana_10/portafolio/.gitkeep"
# Archivos raiz
touch README.md .gitignore
\end{lstlisting}

\begin{tcolorbox}[cajaMejora,title={\faIcon{exclamation-triangle}~No olvides los archivos base}]
Si usas la opción B, debes crear también el contenido de \texttt{README.md} y \texttt{.gitignore} tomando como referencia la plantilla oficial. Un repositorio sin \texttt{.gitignore} corre el riesgo de subir archivos temporales, entornos virtuales o secretos.
\end{tcolorbox}

\section{\iconotexto{user-edit}{Personaliza la plantilla antes del primer commit}}

Antes de subir nada a GitHub, revisa estos puntos:

\begin{enumerate}[leftmargin=1.6em,itemsep=4pt,topsep=2pt]
  \item \textbf{README.md raíz:} reemplaza los marcadores \texttt{[Nombre del integrante N]} por los nombres reales del grupo. Completa la tabla de roles y la tabla de semanas.
  \item \textbf{.gitignore:} confirma que ignora \texttt{.env}, \texttt{.venv/}, \texttt{\_\_pycache\_\_/} y auxiliares de LaTeX. El \texttt{informe.pdf} \textbf{sí} debe entregarse, por eso no se ignora.
  \item \textbf{Sin secretos:} nunca escribas contraseñas, tokens, claves API ni correos privados dentro de archivos versionados. Si necesitas variables locales, usa \texttt{.env} y déjalo ignorado.
  \item \textbf{Sin datos personales sensibles:} no subas cédulas, teléfonos ni información privada innecesaria.
\end{enumerate}

\begin{cajaReto}
Haz una revisión rápida con \texttt{git status} después de \texttt{git add}. Si ves aparecer \texttt{.env}, \texttt{venv/} o archivos grandes inesperados, detente y corrige el \texttt{.gitignore} antes de hacer commit.
\end{cajaReto}

\section{\iconotexto{github}{Crear el repositorio privado en GitHub}}

\begin{enumerate}[leftmargin=1.6em,itemsep=4pt,topsep=2pt]
  \item Entra a \url{https://github.com} e inicia sesión.
  \item Haz clic en \textbf{New repository} (o el botón \texttt{+} \textrightarrow{} \texttt{New repository}).
  \item Configura el repositorio:
  \begin{itemize}[leftmargin=1.4em,itemsep=2pt,topsep=2pt]
    \item \textbf{Repository name:} por ejemplo \texttt{lab2-fisica-grupo-01}.
    \item \textbf{Visibility:} \textbf{Private}. El trabajo del grupo no debe ser público.
    \item \textbf{No} marques "Add a README file", "Add .gitignore" ni "Choose a license" si ya vas a subir la plantilla local; así evitas conflictos en el primer push.
  \end{itemize}
  \item Crea el repositorio y copia la URL HTTPS, que se verá así:
\end{enumerate}

\begin{lstlisting}
https://github.com/<usuario-del-grupo>/lab2-fisica-grupo-01.git
\end{lstlisting}

\begin{tcolorbox}[tarjetaDato,title={\faIcon{user-plus}~Agregar al docente como colaborador}]
En el repositorio de GitHub, ve a \textbf{Settings} \textrightarrow{} \textbf{Collaborators} \textrightarrow{} \textbf{Add people} e invita al docente con permiso de lectura. El docente necesita acceso para revisar las entregas semanales. Confirma con el docente el usuario o correo que debes invitar.
\end{tcolorbox}

\section{\iconotexto{upload}{Conectar tu carpeta local y hacer push}}

Abre una terminal \textbf{dentro} de la carpeta de tu repositorio local (la que contiene \texttt{README.md}, \texttt{.gitignore} y las carpetas \texttt{semana\_01} ... \texttt{semana\_10}).

\subsection{\iconotexto{play-circle}{Comandos para el primer push}}

\begin{lstlisting}
# 1) Inicializar Git (si la carpeta aun no es un repo)
git init

# 2) Ver el estado de los archivos
git status

# 3) Agregar todos los archivos de la plantilla
git add .

# 4) Crear el primer commit
git commit -m "Inicializa estructura del repositorio del grupo"

# 5) Asegurar que la rama principal se llame main
git branch -M main

# 6) Conectar con GitHub (reemplaza la URL por la tuya)
git remote add origin https://github.com/<usuario-del-grupo>/lab2-fisica-grupo-01.git

# 7) Subir la estructura a GitHub
git push -u origin main
\end{lstlisting}

\begin{cajaConcepto}
\texttt{git push -u origin main} sube la rama \texttt{main} y deja configurado el remoto \texttt{origin} como destino por defecto. A partir de ahí, los siguientes envíos pueden hacerse con \texttt{git push}.
\end{cajaConcepto}

\subsection{\iconotexto{key}{Autenticación al hacer push}}

GitHub ya no acepta contraseña de cuenta para operaciones Git por HTTPS. Tienes dos opciones comunes:

\begin{itemize}[leftmargin=1.4em,itemsep=4pt,topsep=2pt]
  \item \textbf{HTTPS + token (PAT):} cuando Git pida contraseña, usa un \emph{Personal Access Token} creado en GitHub (\textbf{Settings} \textrightarrow{} \textbf{Developer settings} \textrightarrow{} \textbf{Personal access tokens}). Guarda el token en un lugar seguro.
  \item \textbf{SSH:} genera una llave SSH, agrégala al agente local y registra la llave pública en GitHub. Luego usa la URL SSH del repositorio:
\end{itemize}

\begin{lstlisting}
# Ejemplo de URL SSH (alternativa a HTTPS)
git remote add origin git@github.com:<usuario-del-grupo>/lab2-fisica-grupo-01.git
\end{lstlisting}

\subsection{\iconotexto{sync-alt}{Subir cambios posteriores}}

Después del primer push, el flujo semanal es simple:

\begin{lstlisting}
git status
git add semana_01
git commit -m "Agrega datos y analisis de la semana 01"
git push
\end{lstlisting}

\begin{cajaRecuerda}
Haz commits pequeños y descriptivos. No esperes al final del semestre para subir todo: el repositorio es la evidencia continua del trabajo del grupo.
\end{cajaRecuerda}

\section{\iconotexto{check-double}{Verificación final}}

Después de hacer \texttt{git push}, entra a tu repositorio en GitHub y revisa:

\begin{tcolorbox}[tarjetaIcono,title={\faIcon{clipboard-check}~Checklist de verificación}]
\begin{itemize}[leftmargin=1.4em,itemsep=3pt,topsep=2pt]
  \item[\faIcon{check-square}] El repositorio es \textbf{privado}.
  \item[\faIcon{check-square}] El docente fue agregado como colaborador.
  \item[\faIcon{check-square}] Existe \texttt{README.md} en la raíz con integrantes y roles.
  \item[\faIcon{check-square}] Existe \texttt{.gitignore} en la raíz.
  \item[\faIcon{check-square}] Están las carpetas \texttt{semana\_01} a \texttt{semana\_10}.
  \item[\faIcon{check-square}] Cada semana tiene \texttt{datos/}, \texttt{analisis/}, \texttt{figuras/} y \texttt{README.md}.
  \item[\faIcon{check-square}] La semana 10 incluye \texttt{portafolio/}.
  \item[\faIcon{check-square}] No aparecen archivos \texttt{.env}, \texttt{venv/} ni credenciales.
\end{itemize}
\end{tcolorbox}

\section{\iconotexto{life-ring}{Problemas comunes y cómo resolverlos}}

\begin{tcolorbox}[cajaMejora,title={\faIcon{bug}~Errores frecuentes}]
\begin{itemize}[leftmargin=1.4em,itemsep=4pt,topsep=2pt]
  \item \textbf{\texttt{remote origin already exists}:} el remoto ya fue configurado. Verifica con \texttt{git remote -v}. Si la URL está mal, corrígela con \texttt{git remote set-url origin <URL-correcta>}.
  \item \textbf{Rechazo por autenticación:} revisa si usaste HTTPS con token o SSH con llave configurada. No uses la contraseña de tu cuenta.
  \item \textbf{El push falla porque el remoto tiene commits:} no debería pasar si creaste el repo vacío. Si pasa, revisa con el docente antes de forzar nada. Evita \texttt{git push --force} en el repositorio del grupo.
  \item \textbf{No veo \texttt{.gitignore} en mi explorador:} es un archivo oculto. En Linux/macOS usa \texttt{ls -a}; en Windows habilita "Elementos ocultos".
  \item \textbf{Subí un archivo que no debía:} elimínalo del índice con \texttt{git rm --cached <archivo>}, agrega la regla a \texttt{.gitignore}, haz commit y push. Si era un secreto, considera el secreto comprometido y cámbialo.
  \item \textbf{Carpetas vacías no aparecen en GitHub:} Git no versiona carpetas vacías. La plantilla usa \texttt{.gitkeep} para conservarlas.
\end{itemize}
\end{tcolorbox}

\begin{tcolorbox}[cajaFortaleza]
Si sigues esta estructura desde la semana 1, el cierre del semestre será mucho más fácil: tendrás datos trazables, scripts reproducibles, figuras regenerables y un historial claro de quién hizo qué.
\end{tcolorbox}

\end{document}
